\chapter{Finite fields}\label{chap:M07.13}
We write $q=p^r$ with $p$ prime and $r\geq1$. Once an algebraic closure of $\F_q$ is fixed, $\F_{q^n}$ denotes its particular subfield of $q^n$ elements, not an arbitrary unidentified copy. Existence and uniqueness of that embedded subfield are proved below.

\KEntry{K07.13.prime_power_cardinality}{Prime-power cardinality}
\begin{theorem}
A finite field has characteristic a prime $p$ and cardinality $p^n$, where $n=[F:\F_p]\geq1$. If $K$ has $q$ elements and $[L:K]=n<\infty$, then $|L|=q^n$.
\end{theorem}
\begin{proof}
The additive order of $1$ is finite, say $p$. If $p=ab$ with $1<a,b<p$, then $(a1)(b1)=0$ contradicts absence of zero divisors and minimality; hence $p$ is prime. The multiples of $1$ form a copy of $\Z/p\Z=\F_p$. A finite field has finite dimension over this subfield: repeatedly selecting an element outside the span must stop because the underlying set is finite. If a basis has $n$ members, its coordinates are arbitrary $n$-tuples from $\F_p$, giving $p^n$ elements. The same coordinate bijection over a $q$-element base gives $q^n$. A field has nonzero dimension, since $1\ne0$.
\end{proof}

\KEntry{K07.13.power_identity}{Power identities}
\begin{theorem}
In a $q$-element field $F$, $a^q=a$ for every $a$, and $a^{q-1}=1$ for $a\ne0$. Moreover
\[X^q-X=\prod_{a\in F}(X-a),\]
and all its roots are simple.
\end{theorem}
\begin{proof}
The group $F^\times$ has $q-1$ elements. Lagrange's theorem gives the nonzero identity; multiplying by $a$ and separately considering $a=0$ gives the first. The displayed monic polynomial has $q$ distinct known roots and degree $q$, so successive root factorization gives the product. Its derivative is $qX^{q-1}-1=-1$ because $p\mid q$, another proof of simplicity.
\end{proof}

\KEntry{K07.13.frobenius}{Absolute and relative Frobenius}
\begin{theorem}
In characteristic $p$, $\phi(x)=x^p$ is an injective field endomorphism, surjective on a finite field. If $K=\F_q$ and $L/K$ is algebraic, $\Phi(x)=x^q$ is a $K$-automorphism of $L$.
\end{theorem}
\begin{proof}
For $0<i<p$, $p$ divides $\binom pi$, so $(x+y)^p=x^p+y^p$. Multiplication and $1$ are preserved, and $(x-y)^p=0$ implies $x=y$. Injectivity is surjectivity on a finite set. Iterate this homomorphism $r$ times to get the $q$-power map; it fixes $K$ by the preceding identity. For any $a\in L$, the field $K(a)$ is finite-dimensional over finite $K$ and hence a finite set. The $q$-power map is a bijection of this field, so some $b\in K(a)$ satisfies $b^q=a$. This proves surjectivity on all of $L$. It does not claim absolute $p$-Frobenius fixes $\F_{p^r}$ when $r>1$.
\end{proof}

\begin{remark}
Relative Frobenius has infinite order on the algebraic closure $\overline{\F_q}$. If a positive power $\Phi^j$ were the identity, every element of that closure would be a root of the nonzero degree-$q^j$ polynomial $X^{q^j}-X$. But Entry~\ref{K07.13.existence} constructs in the closure a subfield with $q^{j+1}>q^j$ elements, contradicting the root bound. That existence proof uses only the Frobenius homomorphism law, not the infinite-order assertion, so this forward reference has no circular dependency.
\end{remark}

\KEntry{K07.13.existence}{Existence of finite fields}
\begin{theorem}
For every prime $p$ and $n\geq1$ a field of $p^n$ elements exists. More generally, in an algebraic closure of a finite field $K$ of cardinality $q$, the roots of $X^{q^n}-X$ form a degree-$n$ extension of $K$ and are its splitting field.
\end{theorem}
\begin{proof}
Take a splitting field of $X^{q^n}-X$. Its derivative is $-1$, so it has exactly $q^n$ distinct roots. They form a subfield: the $q^n$-power endomorphism preserves sums, differences and products, and if $a^{q^n}=a\ne0$, then $(a^{-1})^{q^n}=a^{-1}$. Both $0,1$ are roots. Each $a\in K$ is a root, because iterating $a^q=a$ gives $a^{q^n}=a$. Thus this root set is a field containing $K$, so it is already the field generated by the roots, and equals the splitting field. Its cardinality and the coordinate-count theorem give degree $n$, since $q^d=q^n$ implies $d=n$ for $q>1$. Taking $K=\F_p$ proves the first assertion.
\end{proof}

\KEntry{K07.13.uniqueness}{Uniqueness with the base fixed}
\begin{theorem}
Two fields of order $p^n$ are isomorphic. If $E,L$ are extensions of the same finite field $K$ and have equal finite degree $n$, they are $K$-isomorphic. Inside a fixed algebraic closure of $K$ there is exactly one field of $q^n$ elements containing $K$.
\end{theorem}
\begin{proof}
Every $q^n$-element field over $K$ consists entirely of roots of $X^{q^n}-X$, by the power identity; these elements generate it. Hence it is a splitting field of that polynomial. Splitting-field uniqueness gives the $K$-isomorphism. In a fixed closure the root set itself is specified, so the fields are literally equal. Equal-cardinality abstract fields of characteristic $p$ have their canonical prime-field embeddings, so the same argument over $\F_p$ applies. There is no preferred choice among the possible isomorphisms.
\end{proof}

\KEntry{K07.13.multiplicative_cyclic}{Cyclic unit groups}
\begin{theorem}
Every finite subgroup of a field's multiplicative group is cyclic. Thus $\F_q^\times$ is cyclic of order $q-1$ and has $\varphi(q-1)$ generators.
\end{theorem}
\begin{proof}
Lemma~\ref{lem:finite-multiplicative-cyclic} proves the first assertion using the root bound and a maximal-order element. In a cyclic group $\langle g\rangle$ of order $m$, $g^a$ has order $m/\gcd(a,m)$: $m\mid ak$ is equivalent to $m/\gcd(a,m)\mid k$. Thus it is a generator exactly when the residue class of $a$ is coprime to $m$, and there are $\varphi(m)$ such residues. Apply $m=q-1$. This statement concerns units; the additive group of $\F_{p^n}$ has exponent $p$ and is cyclic only for $n=1$.
\end{proof}

\KEntry{K07.13.primitive_distinction}{Multiplicative generators and field generators}
\begin{proposition}
A generator $a$ of $\F_{q^n}^\times$ generates $\F_{q^n}$ as a field over $\F_q$. A field generator need not generate the multiplicative group.
\end{proposition}
\begin{proof}
The field $\F_q(a)$ contains $0$ and every power of $a$, hence every element. For the converse counterexample, $X^2+1$ has no root in $\F_3$ and is irreducible. In $\F_9=\F_3(i)$ with $i^2=-1$, the element $i$ generates the field but has order $4$, while its unit group has order $8$.
\end{proof}

\KEntry{K07.13.subfield_lattice}{Subfields and divisors}
\begin{theorem}
The subfields of $\F_{p^n}$ are the unique fields $\F_{p^d}$ for positive $d\mid n$. More generally the intermediate fields of $\F_{q^n}/\F_q$ are the unique $\F_{q^d}$ for $d\mid n$, consisting of the roots of $X^{q^d}-X$.
\end{theorem}
\begin{proof}
For an intermediate field $E$, write $d=[E:\F_q]$. The tower law gives $d\mid n$, and $E$ has $q^d$ elements. Uniqueness inside the closure identifies it with that root set. Conversely, if $n=dk$, every root of $X^{q^d}-X$ also satisfies $x^{q^n}=x$ by iterating $k$ times. Thus the already constructed root field of $q^d$ elements is contained in $\F_{q^n}$. Take $q=p$ for all subfields, since every subfield contains the prime field. Positivity is essential; no subfield has just one element.
\end{proof}

\KEntry{K07.13.intersection_compositum}{Intersections and composita}
\begin{theorem}
In the fixed closure of $\F_q$,
\[\F_{q^a}\cap\F_{q^b}=\F_{q^{\gcd(a,b)}},\qquad
\F_{q^a}\F_{q^b}=\F_{q^{\lcm(a,b)}}.\]
Also $\F_{q^a}\subseteq\F_{q^b}$ exactly when $a\mid b$.
\end{theorem}
\begin{proof}
The inclusion assertion follows from the preceding classification. The intersection has degree $d$ dividing both $a,b$, while it contains the degree-$\gcd(a,b)$ field. Divisibility in both directions makes its degree $\gcd(a,b)$, and uniqueness identifies it. Both fields lie in the degree-$\lcm(a,b)$ field; their compositum is therefore finite, of degree $c$ dividing that least common multiple. Since it contains each field, both $a,b$ divide $c$. Hence $c=\lcm(a,b)$, again giving equality of the embedded fields.
\end{proof}

\KEntry{K07.13.galois_group}{Relative finite-field Galois groups}
\begin{theorem}
$L=\F_{q^n}$ is normal and separable over $K=\F_q$, and its $K$-automorphism group is cyclic of order $n$, generated by $\Phi(x)=x^q$. Thus it is a finite Galois extension in the terminology of the next chapter.
\end{theorem}
\begin{proof}
The field is a splitting field, so is normal; its base is perfect, so is separable. The automorphism $\Phi$ has $\Phi^n=1$. If $0<j<n$ and $\Phi^j=1$, every one of the $q^n$ elements would be a root of the nonzero degree-$q^j$ polynomial $X^{q^j}-X$, contradicting the root bound. Thus $\Phi$ has order $n$. There are at most $[L:K]=n$ automorphisms by the embedding count. Its $n$ powers already give that many, so they are all the automorphisms.
\end{proof}

\KEntry{K07.13.embeddings}{Embeddings over a shared finite base}
\begin{theorem}
Let $K=\F_q$ and let finite $K$-extensions $E,L$ have degrees $a,b$. A $K$-embedding $E\to L$ exists if and only if $a\mid b$; when it exists there are exactly $a$ such embeddings.
\end{theorem}
\begin{proof}
An embedding realizes $E$ as an intermediate degree-$a$ field of $L$, so the tower law gives divisibility. Conversely identify $E,L$ over $K$ with the unique fields of degrees $a,b$ in the fixed closure. Divisibility gives inclusion. All $K$-embeddings of $E$ into the closure have the same image, the unique degree-$a$ field, contained in $L$. The finite separable embedding count is $a$, so all $a$ embeddings have codomain $L$ and there are no others. Shared characteristic alone does not replace a specified shared base field.
\end{proof}

\KEntry{K07.13.embedding_compatibility}{Compatible embedding towers}
\begin{proposition}
If $a\mid b\mid c$ are positive, embeddings of the degree-$a$, degree-$b$, and degree-$c$ extensions of $\F_q$ can be chosen compatibly. In particular, for a prescribed $K$-embedding $\sigma:E_a\to E_c$, one can choose $\iota:E_a\to E_b$ and $\tau:E_b\to E_c$ with $\tau\iota=\sigma$.
\end{proposition}
\begin{proof}
Choose any embedding $\iota$, which exists by divisibility. The map $\sigma\iota^{-1}$ on $\iota(E_a)$ extends to an embedding of $E_b$ into an algebraic closure of $E_c$. Its image is the unique degree-$b$ extension of $K$ in that closure and hence is contained in $E_c$, by $b\mid c$. Restrict its codomain to obtain $\tau$. For an entire coherent family, use the unique root fields inside one closure and literal inclusions. Arbitrary independently chosen isomorphisms need not make a triangle commute; compatibility has been constructed here.
\end{proof}

\KEntry{K07.13.frobenius_roots}{Irreducible roots as Frobenius orbits}
\begin{theorem}
If $f\in\F_q[X]$ is irreducible of degree $d$ and $f(a)=0$, its splitting field is $\F_q(a)=\F_{q^d}$. Its distinct roots are $a,a^q,\ldots,a^{q^{d-1}}$, and $d$ is the least positive $j$ with $a^{q^j}=a$.
\end{theorem}
\begin{proof}
The simple-degree theorem gives $[\F_q(a):\F_q]=d$, so this field has $q^d$ elements and $\Phi^d(a)=a$. Every $a^{q^j}$ is a root because Frobenius fixes the coefficients. If $0<j<d$ fixed $a$, it would fix the field generated by $a$, making all $q^d$ elements roots of $X^{q^j}-X$, impossible. Thus the first $d$ iterates are distinct, account for all roots, and already lie in the simple root field. Their product $(X-a)\cdots(X-a^{q^{d-1}})$ is the monic associate of $f$. This proves every assertion, including separability and the splitting-field identity.
\end{proof}

\KEntry{K07.13.universal_factorization}{The universal splitting polynomial}
\begin{theorem}
For $n\geq1$, $X^{q^n}-X$ is the product, each factor once, of all monic irreducibles over $\F_q$ whose degrees divide $n$. Equivalently an irreducible $f$ divides it exactly when $\deg f\mid n$.
\end{theorem}
\begin{proof}
A root $a$ of an irreducible degree-$d$ factor lies in $\F_{q^n}$ exactly when its generated field $\F_{q^d}$ is a subfield, which is equivalent to $d\mid n$. The evaluation-kernel criterion makes this the asserted polynomial divisibility. Factor $X^{q^n}-X$ into monic irreducibles using polynomial unique factorization. Its derivative $-1$ prevents repetitions, and the divisibility criterion identifies exactly the factors that occur. No infinite product is involved: a finite-degree polynomial has finitely many divisors.
\end{proof}

\KEntry{K07.13.irreducibles_exist}{Irreducibles of every positive degree}
\begin{corollary}
For every $n\geq1$ there is a monic irreducible degree-$n$ polynomial over $\F_q$.
\end{corollary}
\begin{proof}
Construct $\F_{q^n}$ and choose a generator of its cyclic multiplicative group. It generates the field, so its monic minimal polynomial has degree $n$ by the power-basis theorem. The assertion is not true over an algebraically closed coefficient field when $n>1$.
\end{proof}

\KEntry{K07.13.irreducible_count}{Counting monic irreducibles}
\begin{theorem}
If $N_q(n)$ is the number of monic irreducibles of degree $n\geq1$, then
\[q^n=\sum_{d\mid n}dN_q(d),\qquad
N_q(n)=\frac1n\sum_{d\mid n}\mu(d)q^{n/d},\]
where $\mu(1)=1$, $\mu(d)=(-1)^s$ for a product of $s$ distinct primes, and $\mu(d)=0$ when a prime square divides $d$.
\end{theorem}
\begin{proof}
There are only $q^n$ monic polynomials of degree $n$, so every counted set is finite. Comparing degrees in the universal factorization gives the first identity. For any positive $m$, expanding $\prod_{\ell\mid m,\ \ell\text{ prime}}(1-1)$ gives $\sum_{d\mid m}\mu(d)=0$ unless $m=1$, when it equals $1$. Put $A(d)=dN_q(d)$. Then
\[\sum_{e\mid n}\mu(e)q^{n/e}
 =\sum_{e\mid n}\mu(e)\sum_{d\mid n/e}A(d)
 =\sum_{d\mid n}A(d)\sum_{e\mid n/d}\mu(e)=A(n).
\]
Divide in $\Q$ by the positive integer $n$. The equality with the integer count proves integrality; no premature natural-number division is used.
\end{proof}

\KEntry{K07.13.irreducibility_test}{An irreducibility test}
\begin{theorem}
For monic $f\in\F_q[X]$ of degree $n\geq1$, irreducibility is equivalent to
\[X^{q^n}\equiv X\pmod f,\qquad
\gcd(f,X^{q^{n/\ell}}-X)=1\quad\text{for every prime }\ell\mid n.
\]
\end{theorem}
\begin{proof}
If $f$ is irreducible, the degree-divisibility criterion gives the first assertion and says it cannot divide a polynomial with exponent $n/\ell<n$, giving the gcd condition. Conversely the first assertion says $f\mid X^{q^n}-X$. Hence $f$ is squarefree and each irreducible factor has degree $d\mid n$. If $d<n$, some prime $\ell$ divides $n/d$, so $d\mid n/\ell$. That factor would divide $X^{q^{n/\ell}}-X$, contradicting the gcd condition. All factors therefore have degree $n$; the total degree forces exactly one. For $n=1$ the list of prime conditions is empty and the same argument applies.
\end{proof}

\KEntry{K07.13.coordinates_arithmetic}{Polynomial coordinates and arithmetic}
\begin{theorem}
For monic irreducible $f\in\F_q[X]$ of degree $n$, the field $\F_q[X]/(f)$ has $q^n$ elements, uniquely represented by polynomials of degree less than $n$. Addition and multiplication are polynomial operations followed by remainder; a nonzero residue $g$ has inverse obtained from $uf+vg=1$.
\end{theorem}
\begin{proof}
The power basis gives the coordinate bijection $\F_q^n\to\F_q[X]/(f)$, hence the cardinality. Replacing a representative by a multiple of $f$ changes neither the sum nor product class, by the quotient-ring laws. Euclidean division gives the unique reduced representative. For a nonzero class, $f\nmid g$, so irreducibility gives gcd $1$ and Bezout gives the displayed inverse. A nonzero polynomial can still represent zero when it is a multiple of $f$; it is the residue class that must be nonzero.
\end{proof}

\KEntry{K07.13.explicit_F4}{The four-element field}
\begin{example}
$X^2+X+1$ has values $1$ at both $0,1\in\F_2$, so is irreducible. In its quotient write $u=X+(X^2+X+1)$. Then
\[\F_4=\{0,1,u,u+1\},\qquad u^2=u+1,\quad u(u+1)=1,\quad (u+1)^2=u.
\]
\end{example}
\begin{proof}
Degree-two coordinate uniqueness distinguishes these four elements. The defining equation gives $u^2=u+1$ because negatives equal positives in characteristic $2$. It follows that $u(u+1)=u^2+u=1$ and $(u+1)^2=u^2+1=u$. Thus $u^3=1$ while $u\ne1$, so $u$ has order $3$ and generates all nonzero elements. Addition is coefficientwise modulo $2$; these identities determine every product.
\end{proof}

\KEntry{K07.13.primitive_polynomials}{Primitive polynomials}
\begin{definition}
A monic degree-$n$ polynomial, $n\geq1$, over $\F_q$ is primitive when it is irreducible and one of its roots generates $\F_{q^n}^\times$.
\end{definition}
\begin{proposition}
The definition is independent of which root is chosen.
\end{proposition}
\begin{proof}
All roots of an irreducible polynomial are $a^{q^j}$. If $a$ has multiplicative order $m$, then $m\mid q^n-1$, so $\gcd(q,m)=1$. The order-of-a-power formula gives $\operatorname{ord}(a^{q^j})=m$. All roots therefore generate the unit group together or fail together. Irreducibility alone is weaker: $X^2+1$ over $\F_3$ has roots of order $4$, not $8$. This use of ``primitive'' is unrelated to coefficient content $1$ in Gauss's lemma.
\end{proof}

\KEntry{K07.13.polynomial_functions}{Polynomial functions}
\begin{theorem}
$X^q-X$ is a nonzero polynomial inducing the zero function on $\F_q$. Every function $h:\F_q\to\F_q$ has a unique polynomial representative of degree less than $q$.
\end{theorem}
\begin{proof}
Nonzeroness follows from leading coefficient $1$ and degree $q>1$; zero evaluation is the power identity. For the second assertion define
\[P_h(X)=\sum_{a\in\F_q}h(a)\prod_{b\in\F_q,\ b\ne a}\frac{X-b}{a-b}.
\]
Every denominator is nonzero. Evaluation at a node $c$ leaves just the $a=c$ term, giving $P_h(c)=h(c)$, and every summand has degree at most $q-1$. If two degree-below-$q$ polynomials agree at all $q$ nodes, their difference either is zero or violates the root bound. This proves uniqueness without identifying a formal polynomial with its evaluation function.
\end{proof}

\KEntry{K07.13.trace_norm}{Trace and norm}
\begin{theorem}
For $L=\F_{q^n}$ over $K=\F_q$, define $\Tr_{L/K}(x)$ and $\Nm_{L/K}(x)$ as the trace and determinant of multiplication by $x$ on the $K$-vector space $L$. Then
\[\Tr_{L/K}(x)=\sum_{j=0}^{n-1}x^{q^j},\qquad
\Nm_{L/K}(x)=\prod_{j=0}^{n-1}x^{q^j}.
\]
For $x\ne0$, the latter is $x^{(q^n-1)/(q-1)}$. Trace is $K$-linear and surjective, and norm maps $L^\times$ onto $K^\times$.
\end{theorem}
\begin{proof}
Choose $L=K(a)$ with power basis and the $n$ embeddings $\sigma_j=\Phi^j$ into an algebraic closure. If $M_x$ is the multiplication matrix, the row $(1,\sigma_j(a),\ldots,\sigma_j(a)^{n-1})$ satisfies $v_j M_x=\sigma_j(x)v_j$, by evaluating the coordinate identity for multiplication by $x$. The matrix with rows $v_j$ is invertible: a vector in its kernel gives a degree-below-$n$ polynomial vanishing at $n$ distinct conjugates, so is zero. Thus $M_x$ becomes diagonal with entries $\sigma_j(x)$. D05 trace/determinant invariance under change of basis gives the two formulas; the geometric sum of the exponents gives the power formula. Linearity of trace also follows directly from the sum formula and Frobenius fixing $K$.

The polynomial $T(X)=X+X^q+\cdots+X^{q^{n-1}}$ is nonzero of degree $q^{n-1}<q^n$, so cannot vanish on all $L$. Hence trace is a nonzero $K$-linear map into the one-dimensional space $K$ and is onto: scale one value different from zero to any prescribed value. For norm, choose a generator $g$ of $L^\times$ of order $q^n-1$. Its norm has order $q-1$ by the order-of-a-power formula, so generates $K^\times$. This proves surjectivity there (and zero maps to zero). Although $\Tr(1)=n$ can vanish in characteristic $p\mid n$, this does not make the trace map zero.
\end{proof}
